\scanpage{24}
\begin{center}-\mspage{0.16 bis}-\end{center}
\[
\operatorname{tr.d.}\ \unsure{\kappa|\mathcal K}(y)/k+\sigma(y)=\dim G.
\]
Now if the action were not closed, at least over some $k$ there would be two points $x_1,x_2\in X_k$ such that $O(x_2)\subset\overline{O(x_1)}-O(x_1)$. Let $y_i$ be the generic point of $O(x_i)$. Since $x_2$ is a regular point of the action, one finds that $\sigma(x_1)=\sigma(x_2)$, hence
\[
\operatorname{tr.d.}\ \unsure{\kappa|\mathcal K}(y_1)/k=\operatorname{tr.d.}\ \unsure{\kappa|\mathcal K}(y_2)/k<\infty.
\]
But $y_2$ is a specialization of $y_1$, hence this is absurd.

Concerning restrictive conditions on the action of $G$, a notion intermediate between the closedness of one orbit and the properness of the whole action is sometimes convenient:

\uline{lemma} \uline{0.3}: Let $S=\operatorname{Spec}(k)$, $k$ an algebraically closed field, and assume $X$ and $G$ are algebraic pre-schemes over $k$. Let $\sigma$ be an action of $G$ on $X$, and let $x\in X_k$. Then $\psi_x$ is proper if and only if $O(x)$ is closed in $X$ and $S(x)$ is proper over $k$.

\uline{Proof}: If $\psi_x$ is proper, then its image $O(x)$ must be closed, and its fibre over $x$ - which is $S(x)$ - must be proper over $k$.

On the other hand, assume that $O(x)$ is closed and $S(x)$ is proper over $k$. To prove that $\psi_x$ is proper, we may as well replace $G$
\mspage{23}
